%==============================================================================
% Chapter 4 — System Design
% الهيكل المعتمد (أكتوبر 2026) — كل قسم في ملف مستقل لتوزيع العمل على الملاك
% وفق حزمة الفصل الرابع في الخطة الزمنية (مهمات جانت ٢.٢):
%   ٢.٢.١ المعمارية العامة متعددة الوكلاء (AI)      → sec1 + sec2 + sec3
%   ٢.٢.٢ قواعد البيانات العلائقية والمتجهية (BE)   → sec4
%   ٢.٢.٣ الواجهات ولوحة التحكم الإدارية (FE)       → sec7
%   ٢.٢.٤ مخططات التدفق والتتابع والكيانات (TM)     → موزعة داخل الأقسام
%   ٢.٢.٥ واجهات برمجة التطبيقات ومواصفاتها (BE)    → sec6
%   ٢.٢.٦ إطار الأمان والصلاحيات (BE)               → sec5
%   ٢.٢.٧ مراجعة الفصل واعتماده (TM)                → sec8 + sec9
%
% قواعد ملزمة للكتابة:
%   (1) المخططات لغة توثيق داخل الأقسام، لا قسماً مستقلاً — مهمة ٢.٢.٤ تُنجز
%       داخل أقسامها: المكونات في 4.1، ER في 4.4، السلكية في 4.7...
%   (2) لا إعادة رسم لشكلي الفصل الثالث (fig:agent-state-graph و
%       fig:agent-sequence) — يُحال إليهما بـ \ref{}؛ هذا الفصل يتعمق في
%       «كيف صُمم» لا «كيف يتفاعل»
%   (3) كل عنصر تصميم يُربط بمتطلب وظيفي (FR-01..FR-29) أو حالة استخدام
%       (UC-01..UC-07)، وتُجمع الروابط في مصفوفة القسم 4.8
%   (4) الأشكال الجديدة TikZ تُحفظ في figures/tikz/ وتُدرج بـ \input،
%       والجداول تثبت بموضعها [H]
%==============================================================================
\chapter{تصميم النظام}
\label{ch:design}
\markboth{الفصل الرابع: تصميم النظام}{الفصل الرابع: تصميم النظام}

%------------------------------------------------------------------------------
% 4.1 المعمارية العامة للنظام
% المالك: AI (مهمة جانت ٢.٢.١ — التسليم 27/10/2026)
% الحجم المستهدف: 3-4 صفحات
% الغرض: الانتقال من التحليل (الفصل الثالث) إلى القرارات التصميمية الكلية:
%   لماذا هذه البنية تحديداً، وكيف تستجيب للمتطلبات غير الوظيفية.
%
% عناصر مطلوبة:
%   [ ] 4.1.1 وصف النمط المعماري: وكلاء لغويون منسقون برسم بياني حالة
%       (LangGraph) فوق بنية ثلاثية الطبقات (واجهة Next.js / خدمات FastAPI /
%       تخزين مزدوج PostgreSQL + Qdrant)، مع دور المنسق المركزي (Supervisor)
%   [ ] 4.1.2 شكل مخطط المكونات UML (fig:component-diagram) — TikZ جديد في
%       figures/tikz/component-diagram.tex: المكونات = واجهة المحادثة، خادم
%       API، الوكلاء الخمسة، طبقة العروض الذكية، قاعدة ERP (قراءة فقط)،
%       Qdrant، قاعدة بيانات النظام، وحدة التدقيق — مع واجهات موفرة/مطلوبة.
%       يمتاز عن الشكل المفاهيمي في الفصل الثاني بأنه تصميمي بمستوى
%       المكونات ومنافذها
%   [ ] 4.1.3 جدول القرارات المعمارية (tab:arch-decisions): القرار / البديل
%       المرفوض / المبرر / المتطلب المخدوم — 6-8 قرارات (لماذا LangGraph لا
%       CrewAI أو AutoGen؟ لماذا قراءة فقط من ERP؟ لماذا Qdrant لا pgvector؟
%       لماذا نشر Docker Compose؟ ...)
%
% تتبع: يستجيب للمتطلبات غير الوظيفية (الأداء، القابلية للتوسع، الخصوصية،
%   قابلية الصيانة) بأرقامها من جدول المتطلبات غير الوظيفية في الفصل الثالث
%------------------------------------------------------------------------------
\section{المعمارية العامة للنظام}
\label{sec:architecture}
~

\subsection{النمط المعماري}
\label{subsec:arch-pattern}
~

\subsection{مخطط المكونات}
\label{subsec:arch-components}
~

\subsection{القرارات المعمارية الرئيسة}
\label{subsec:arch-decisions}
~

%------------------------------------------------------------------------------
% 4.2 تصميم منظومة الوكلاء
% المالك: AI (امتداد مهمة ٢.٢.١ — قد يستدعي إعادة توزيع داخل الفريق)
% الحجم المستهدف: 3-4 صفحات
% الغرض: التصميم التنفيذي لما نمذجناه تحليلياً في القسم 3.5 — مستوى العقود
%   والأنواع وقواعد التوجيه، لا مستوى الوصف الوظيفي.
%
% تحذيرات:
%   - لا إعادة رسم للشكل 3-7 (fig:agent-state-graph) ولا الشكل 3-8
%     (fig:agent-sequence): أحِل إليهما بـ \ref{} وابنِ فوقهما
%   - لا تكرار لجدول الحقول التحليلي في 3.5.2: هنا المستوى تصميمي — الأنواع
%     (TypedDict)، القيم الابتدائية، قواعد التحقق، ودوال الاختزال (Reducers)
%
% عناصر مطلوبة:
%   [ ] 4.2.1 طوبولوجيا الرسم البياني ومواصفات العقود (tab:agent-contracts):
%       لكل عقدة — التوقيع (مدخل/مخرج)، السلوك عند الفشل، المهلة القصوى،
%       وسياسة إعادة المحاولة
%   [ ] 4.2.2 تصميم الحالة المشتركة: بنية AgentState التنفيذية وعلاقتها بسجل
%       الجلسة المعمّر في قاعدة بيانات النظام (ما يبقى في الذاكرة وما يُثبَّت)
%   [ ] 4.2.3 منطق الحواف المشروطة: جدول قواعد التوجيه (tab:edge-rules) —
%       الحالة / الشرط / الوجهة التالية؛ يشمل حلقة التصحيح الذاتي لوكيل SQL
%       (حد أقصى 3 محاولات مع تغذية Few-Shot من رسالة الخطأ) — اربطها بمتطلب
%       المتانة المعني من جدول المتطلبات الوظيفية
%   [ ] 4.2.4 تصميم الأدوات (tab:tool-schemas): منفذ SQL للقراءة فقط، مسترجع
%       Qdrant، مصنف النوايا، كاشف القيم الشاذة، مولد التقارير — لكل أداة:
%       المدخل / المخرج / القيود الأمنية
%
% مرجع الدقة: الرسم البياني والأدوات الفعلية في مستودع Bayan
%   (backend/src/agents) — صمم من الكود الحقيقي لا من التصور
%------------------------------------------------------------------------------
\section{تصميم منظومة الوكلاء}
\label{sec:agents-design}
~

\subsection{طوبولوجيا الرسم البياني ومواصفات العقود}
\label{subsec:agents-topology}
~

\subsection{تصميم الحالة المشتركة}
\label{subsec:agents-state}
~

\subsection{منطق الحواف المشروطة وقواعد التوجيه}
\label{subsec:agents-edges}
~

\subsection{تصميم الأدوات}
\label{subsec:agents-tools}
~

%------------------------------------------------------------------------------
% 4.3 تصميم طبقة الرؤى الذكية (Smart Views)
% المالك: AI بالتعاون مع BE
% الحجم المستهدف: 2-3 صفحات
% الغرض: توثيق المساهمة التقنية المميزة للمشروع تصميمياً — الطبقة الوسيطة
%   بين مخطط ERP الخام ووكيل SQL، والأساس الذي يقوم عليه الأمان بالتصنيف
%   الرباعي للجداول.
%
% عناصر مطلوبة:
%   [ ] 4.3.1 بنية الطبقة: تعريف «العرض الذكي» (استعلام SQL معياري + توصيف
%       النوايا العربية المقصودة + تسميات الأعمال للمخرجات) وعلاقته بالجداول
%       المصدر، مع مثال تصميمي واحد كامل (عرض الملف الطلابي
%       v_student_profile من مستودع Bayan): الجداول المصدر، المعاملات،
%       النوايا المرتبطة، الدور المسموح
%   [ ] 4.3.2 التصنيف الأمني رباعي المستويات (tab:sv-security-tiers):
%       عام / داخلي / مقيّد بالدور / سري — لكل مستوى: من يراه، وكيف يقيّد
%       توليد SQL (رفض أم تعويض بأعم؟)، وأثره في سجل التدقيق
%   [ ] 4.3.3 دورة حياة العرض الذكي: التعريف ← المراجعة ← الإصدار ←
%       الاكتشاف التلقائي أثناء التهيئة الذكية (يرتبط بحالة الاستخدام UC-07)،
%       وكيف يتعامل النظام مع تعديل مخطط ERP المصدر (إبطال العرض المعطوب)
%
% تتبع: يخدم مباشرة حالات الاستخدام UC-01 وUC-05 وUC-07 ومتطلبات الاستعلام
%   العربي وربط مصادر البيانات بأرقامها من جدول المتطلبات الوظيفية
%------------------------------------------------------------------------------
\section{تصميم طبقة الرؤى الذكية (\EN{Smart Views})}
\label{sec:smart-views-design}
~

\subsection{بنية الطبقة وعلاقتها بمخطط \EN{ERP}}
\label{subsec:sv-structure}
~

\subsection{التصنيف الأمني رباعي المستويات للجداول}
\label{subsec:sv-security-tiers}
~

\subsection{دورة حياة العرض الذكي}
\label{subsec:sv-lifecycle}
~

%------------------------------------------------------------------------------
% 4.4 تصميم قواعد البيانات والتخزين
% المالك: BE (مهمة جانت ٢.٢.٢ — التسليم 31/10/2026)
% الحجم المستهدف: 3-4 صفحات
%
% عناصر مطلوبة:
%   [ ] 4.4.1 قاعدة بيانات ERP (Odoo Education): فلسفة «القراءة فقط عبر
%       العروض الذكية» — لا تعديل لمخطط Odoo إطلاقاً؛ الجداول الأساسية
%       المستفادة (الطلاب، البرامج، الكليات، الرسوم، النتائج...) ومستخدم
%       الربط المخصص بأدنى صلاحية (SELECT على العروض فقط)
%   [ ] 4.4.2 قاعدة بيانات النظام الداخلية + مخطط ER (fig:bayan-er) — TikZ
%       جديد في figures/tikz/bayan-er.tex: المستخدمون، الأدوار، الجلسات،
%       سجل التدقيق، تعريفات العروض الذكية، فهرس المستندات — بالعلاقات
%       والمفاتيح. راجع backend/prisma/schema.prisma في مستودع Bayan للدقة
%   [ ] 4.4.3 مخزن المتجهات Qdrant: المجموعات (Collections)، بنية الحمولة
%       (Payload: المصدر، رقم الصفحة، التصنيف الأمني للوثيقة)، نموذج التضمين
%       multilingual-e5-large وأبعاد المتجهات، واستراتيجية الفهرسة وإعادة
%       البناء عند تحديث المستندات
%
% تحذيرات:
%   - مخطط ER هنا للقاعدة الداخلية الجديدة فقط؛ مخطط Odoo جاهز ومستقر
%     ويُكتفى بجرد الجداول المستخدمة منه دون إعادة رسمه بالكامل
%   - هذا القسم هو موضع «مخططات الكيانات» التي تعد بها مهمة جانت ٢.٢.٤
%------------------------------------------------------------------------------
\section{تصميم قواعد البيانات والتخزين}
\label{sec:db-design}
~

\subsection{قاعدة بيانات \EN{ERP} بيئة التطبيق}
\label{subsec:erp-db}
~

\subsection{قاعدة البيانات الداخلية ومخطط \EN{ER}}
\label{subsec:bayan-db}
~

\subsection{مخزن المتجهات \EN{Qdrant}}
\label{subsec:vector-store}
~

%------------------------------------------------------------------------------
% 4.5 تصميم الأمان والصلاحيات
% المالك: BE (مهمة جانت ٢.٢.٦ — 09/11 إلى 12/11/2026)
% الحجم المستهدف: 2-3 صفحات
% الغرض: ترجمة سردية الخصوصية في الأطروحة إلى آليات تصميمية قابلة للفحص —
%   اتساقاً مع خطة الاستجابة للمخاطر R-09 بصيغتها المعدلة.
%
% عناصر مطلوبة:
%   [ ] 4.5.1 مصفوفة التحكم المبني على الأدوار RBAC (tab:rbac-matrix):
%       الفاعلون الأربعة (مسؤول المنصة، مسؤول تقنية المعلومات، المدير،
%       الموظف) × حالات الاستخدام السبع — بمستويات (تنفيذ / اطلاع / محجوب)،
%       مع الإشارة إلى علاقة التعميم بين الأدوار الثلاثة المؤسسية
%   [ ] 4.5.2 إدارة الجلسات والرموز: JWT قصيرة الصلاحية، تدفق المصادقة
%       (الدخول / التجديد / الإبطال)، وتحديد معدل المحاولات (Rate Limiting)
%   [ ] 4.5.3 سجل التدقيق وتتبع الوكلاء: ماذا يُسجل لكل طلب (من / متى /
%       ماذا / أي وكيل تدخل / أي عرض ذكي استُخدم / كم استغرق) وكيفية
%       استعلامه وتصفية نتائجه (حالة الاستخدام UC-06)
%
% تحذيرات:
%   - ممنوع إدخال المصادقة الثنائية MFA/2FA في التصميم — الأطروحة تنص على
%     عدم دعمها في هذه النسخة (اتساقاً مع إصلاح تناقض R-09 في الفصل الثالث)
%   - اربط كل آلية بمخاطرها في سجل المخاطر (R-xx) لتظهر السلسلة كاملة:
%     خاطر ← استجابة ← تصميم
%------------------------------------------------------------------------------
\section{تصميم الأمان والصلاحيات}
\label{sec:security-design}
~

\subsection{مصفوفة التحكم المبني على الأدوار (\EN{RBAC})}
\label{subsec:rbac-matrix}
~

\subsection{إدارة الجلسات والرموز}
\label{subsec:auth-session}
~

\subsection{سجل التدقيق وتتبع الوكلاء}
\label{subsec:audit-design}
~

%------------------------------------------------------------------------------
% 4.6 تصميم واجهات برمجة التطبيقات (APIs)
% المالك: BE (مهمة جانت ٢.٢.٥ — 06/11 إلى 09/11/2026)
% الحجم المستهدف: 2-3 صفحات
%
% عناصر مطلوبة:
%   [ ] 4.6.1 الأسلوب المعتمد: REST فوق HTTP/JSON + بث الأحداث من الخادم
%       (SSE) لمسار المحادثة لعرض الإجابة تدريجياً — وهو الجوهر التقني
%       لتجربة الاستجابة السريعة الموعودة في المتطلبات غير الوظيفية
%   [ ] 4.6.2 كتالوج نقاط النهاية (tab:api-endpoints): الطريقة / المسار /
%       الغرض / الدور المسموح — نحو 8-12 نقطة تغطي: المصادقة، المحادثة
%       (والبث)، التقارير والتصدير، إدارة المستخدمين والأدوار، العروض
%       الذكية، سجل التدقيق، التهيئة الذكية، فحص الصحة
%   [ ] 4.6.3 معالجة الأخطاء: صيغة الاستجابة الموحدة (رمز خطأ داخلي ثابت،
%       رسالة عربية للمستخدم، تفاصيل تشخيصية للمطورين) + رموز حالات HTTP
%       المعتمدة لكل عائلة أخطاء (صلاحيات / تحقق من المدخلات / حدود المعدل /
%       أخطاء الخادم)
%
% مرجع الدقة: مسارات API الفعلية في مستودع Bayan (backend/src/api) — ابنِ
%   الكتالوج من الكود الموجود لا من التصور
%------------------------------------------------------------------------------
\section{تصميم واجهات برمجة التطبيقات (\EN{APIs})}
\label{sec:api-design}
~

\subsection{الأسلوب المعتمد}
\label{subsec:api-style}
~

\subsection{كتالوج نقاط النهاية}
\label{subsec:api-catalog}
~

\subsection{معالجة الأخطاء وصيغ الاستجابة الموحدة}
\label{subsec:api-errors}
~

%------------------------------------------------------------------------------
% 4.7 تصميم واجهة المستخدم وتجربتها
% المالك: FE (مهمة جانت ٢.٢.٣ — التسليم 04/11/2026)
% الحجم المستهدف: 2-3 صفحات
%
% عناصر مطلوبة:
%   [ ] 4.7.1 تدفق الحوار العربي (RTL): أنماط المحادثة المعتمدة — اقتراحات
%       الأسئلة الافتتاحية، عرض الاستعلام SQL المنفذ للمستخدم المتقدم،
%       إظهار التوثيق بالمصادر وأرقام الصفحات في إجابات الاسترجاع، أزرار
%       تصدير التقرير، ومعالجة حالات الانتظار والأخطاء بلغة عربية واضحة
%   [ ] 4.7.2 المخططات السلكية للشاشات الرئيسة (fig:wireframes-chat و
%       fig:wireframes-admin) — TikZ في figures/tikz/: شاشة المحادثة الرئيسة
%       + لوحة الإدارة (المستخدمون والأدوار / سجل التدقيق / العروض الذكية)،
%       بأسلوب صناديق رمادية بسيطة (Wireframes) لا رسوم مرئية مزخرفة
%
% ملاحظة: التصميم المرئي النهائي (الهوية والألوان والخطوط) موضعه الطبيعي
%   الفصل الخامس عند عرض التنفيذ — هنا مستوى الهيكل وتدفق المستخدم فقط
%------------------------------------------------------------------------------
\section{تصميم واجهة المستخدم وتجربتها}
\label{sec:ui-design}
~

\subsection{تدفق الحوار العربي وأنماط المحادثة}
\label{subsec:ui-conversation}
~

\subsection{المخططات السلكية للشاشات الرئيسة}
\label{subsec:ui-wireframes}
~

%------------------------------------------------------------------------------
% 4.8 مصفوفة تتبع التصميم إلى المتطلبات
% المالك: TM (مهمة جانت ٢.٢.٧ — المراجعة والاعتماد 12/11 إلى 14/11/2026)
% الغرض: القفل النهائي للفصل — إثبات قابل للفحص أن كل متطلب وظيفي
%   (FR-01..FR-29) وكل حالة استخدام (UC-01..UC-07) لهما موضع تصميم صريح
%   في هذا الفصل، فلا ينتقل شيء إلى التنفيذ بلا أساس تصميمي.
%
% عناصر مطلوبة:
%   [ ] الجدول المرجعي (tab:design-traceability): المتطلب / حالة الاستخدام ←
%       عنصر التصميم المخدوم (رقم القسم أو الجدول أو الشكل) ← ملاحظة —
%       29 صفاً للمتطلبات و7 صفوف لحالات الاستخدام (يجوز تقسيمه جدولين
%       للمتطلبات الوظيفية وحالات الاستخدام)
%   [ ] فقرة تحليلية قصيرة: قراءة المصفوفة — أي متطلبات تخدمها عناصر
%       تصميم متعددة (قوة)، وأيها يخدمها عنصر واحد (نقطة انتباه في
%       الاختبار)، وأي متطلب بلا موضع تصميم (فجوة يجب سدها قبل الفصل
%       الخامس — إن وُجد)
%
% أسلوب الجدول: نفس أسلوب «جدول الملخص المرجعي السريع» الذي نجح في
%   حالات الاستخدام (3.4.4) — أعمدة قصيرة وصف لكل عنصر
%------------------------------------------------------------------------------
\section{مصفوفة تتبع التصميم إلى المتطلبات}
\label{sec:design-traceability}
~

%------------------------------------------------------------------------------
% 4.9 ملخص الفصل
% المالك: TM (ضمن مهمة ٢.٢.٧)
% الغرض: تكثيف القرارات التصميمية (لا سرد العناوين) وربطها بمخرجات الفصل
%   الثالث وتمهيد الفصل الخامس الذي ينفذ النظام على هذه الأسس.
%
% عناصر مطلوبة:
%   [ ] قائمة النتائج التصميمية: 5-7 نقاط جوهرية (المعمارية المعتمدة
%       وقراراتها، عقود الوكلاء وقواعد التوجيه، طبقة العروض الذكية
%       بتصنيفها الرباعي، مخطط ER الداخلي، مصفوفة RBAC، كتالوج API،
%       المصفوفة التتبعية كدليل اكتمال)
%   [ ] جملة الانتقال إلى الفصل الخامس: بيئة التنفيذ والتقنيات الفعلية
%       التي تجسد هذا التصميم
%------------------------------------------------------------------------------
\section{ملخص الفصل}
\label{sec:design-summary}
~

